\section{Appendix U -- Publication-Grade Scientific Quality Standard}
% R005_APPENDIX_READER_FRAME
Reader frame. This appendix supports the main manuscript by explaining source role, evidence class, audit route, and limitation boundary. It is not a detached raw dump.


\label{sec:oc133-public-scientific-quality-standard}

This appendix states the public scientific quality standard applied to OC Core 1.3.3. The enforcement machinery remains in machine-readable release criteria and audit reports; the text here gives the reader-facing method. A release does not pass because it merely avoids forbidden phrases. It passes only when it positively performs the work expected of a scientific monograph, article, methods companion, reviewer map, archive record, and evidence package.


The standard is intentionally stricter than ordinary venue compliance. It combines proof discipline, empirical discipline, prior-art discipline, didactic discipline, editorial discipline, ontology discipline, archive discipline, and post-release accountability. The result is a must-have standard: every public text must answer what it is, who it is for, what it claims, what evidence supports it, how the reader should inspect it, what would falsify it, and what remains outside its present authority.


\subsection{1. Audience, task, and reader contract}

A scientific document has to begin by identifying who it is for, what problem it addresses, what it asks the reader to learn, and what kind of judgment the reader is expected to make. For OC Core this means that the formal-methods reader, the domain scientist, the journal editor, and the technically literate institutional reader must all be able to locate their route through the same text without guessing which artifact carries the scientific argument.


The positive criterion asks whether the document states its audience, its document role, its reader contract, its recommended order, its claim boundary, and its evidence route. A document that contains correct material but never tells the reader why the material is there still fails the publication-grade standard.


Before writing, the editorial system treats audience, task, and reader contract as a design requirement rather than as a late-stage defect class. The outline has to allocate a place for the requirement, the corpus register has to identify the supporting sources, and the claim boundary has to name the evidence class that will carry the requirement. This prevents the manuscript from growing by accidental appendix accretion and forces the scientific argument to declare why the material belongs where it appears.


During manuscript integration, audience, task, and reader contract is checked against three questions. First, does the section teach the reader something that is necessary for judging the released scientific object? Second, does the section cite or explain the artifact layer that supports the statement? Third, does the section give a hostile reviewer a concrete way to challenge the statement? If one of those questions cannot be answered from the text, the section is not ready for release.


At acceptance time, audience, task, and reader contract must leave quantitative traces. The checker records anchors, section counts, evidence references, reader-question coverage, current-science inclusion, repeated boilerplate ratio, stale-identity hits, and public-surface consistency. The exact numbers are not the scientific claim; they are the guardrail that keeps editorial judgment from becoming a vague impression after the fact.


\begin{itemize}[leftmargin=1.8em]

\item \textbf{Positive obligation:} the manuscript must perform this role explicitly, not merely avoid a corresponding defect

\item \textbf{Quantitative witness:} the criterion records anchors, counts, coverage booleans, hashes, or repeated-run stability rather than relying on editorial impression alone

\item \textbf{Repair route:} a failure routes to manuscript integration, scientific editorial, proof repair, empirical replay, comparator repair, public metadata repair, or post-release incident handling according to the affected object class

\end{itemize}

\subsection{2. Current science integration}

The release must integrate the current model state rather than merely attach a new folder to an older manuscript. The model state includes the typed foundation, promoted claim register, theorem inventory, proof sheets, Lean subset, finite semantic checks, target-blind or held-out replay rows, prior-art comparison, phenomenon coverage, reviewer closure, and journal package map. Each surface must either appear in the manuscript as prose, be cited as evidence, or be explicitly excluded with a reason.


The positive criterion therefore measures coverage. It is not enough that a file exists in the repository. The text must tell the reader what the artifact contributes, what it cannot prove, and how a reviewer can check it. Hidden current science is treated as missing current science until it is represented by a public explanation or by a referenced evidence object.


Before writing, the editorial system treats current science integration as a design requirement rather than as a late-stage defect class. The outline has to allocate a place for the requirement, the corpus register has to identify the supporting sources, and the claim boundary has to name the evidence class that will carry the requirement. This prevents the manuscript from growing by accidental appendix accretion and forces the scientific argument to declare why the material belongs where it appears.


During manuscript integration, current science integration is checked against three questions. First, does the section teach the reader something that is necessary for judging the released scientific object? Second, does the section cite or explain the artifact layer that supports the statement? Third, does the section give a hostile reviewer a concrete way to challenge the statement? If one of those questions cannot be answered from the text, the section is not ready for release.


At acceptance time, current science integration must leave quantitative traces. The checker records anchors, section counts, evidence references, reader-question coverage, current-science inclusion, repeated boilerplate ratio, stale-identity hits, and public-surface consistency. The exact numbers are not the scientific claim; they are the guardrail that keeps editorial judgment from becoming a vague impression after the fact.


\begin{itemize}[leftmargin=1.8em]

\item \textbf{Positive obligation:} the manuscript must perform this role explicitly, not merely avoid a corresponding defect

\item \textbf{Quantitative witness:} the criterion records anchors, counts, coverage booleans, hashes, or repeated-run stability rather than relying on editorial impression alone

\item \textbf{Repair route:} a failure routes to manuscript integration, scientific editorial, proof repair, empirical replay, comparator repair, public metadata repair, or post-release incident handling according to the affected object class

\end{itemize}

\subsection{3. Claim-to-evidence binding}

Every scientific assertion must carry an evidence class. Some assertions are definitional, some are theorem-bound, some are mechanized in a selected Lean subset, some are executable finite witnesses, some are replay rows, some are comparator statements, and some are only non-promoted research obligations. The manuscript must not blur those classes. A statement supported by finite examples cannot be promoted as a full empirical law; a comparator row cannot become proof of scoped novelty.


The positive criterion checks that claims mention assumptions, evidence routes, proof or replay artifacts, comparator context, limits, and reopening conditions. The safest release is not the one that says little; it is the one that states ambitious claims only at the level of support actually carried by the artifact layer.


Before writing, the editorial system treats claim-to-evidence binding as a design requirement rather than as a late-stage defect class. The outline has to allocate a place for the requirement, the corpus register has to identify the supporting sources, and the claim boundary has to name the evidence class that will carry the requirement. This prevents the manuscript from growing by accidental appendix accretion and forces the scientific argument to declare why the material belongs where it appears.


During manuscript integration, claim-to-evidence binding is checked against three questions. First, does the section teach the reader something that is necessary for judging the released scientific object? Second, does the section cite or explain the artifact layer that supports the statement? Third, does the section give a hostile reviewer a concrete way to challenge the statement? If one of those questions cannot be answered from the text, the section is not ready for release.


At acceptance time, claim-to-evidence binding must leave quantitative traces. The checker records anchors, section counts, evidence references, reader-question coverage, current-science inclusion, repeated boilerplate ratio, stale-identity hits, and public-surface consistency. The exact numbers are not the scientific claim; they are the guardrail that keeps editorial judgment from becoming a vague impression after the fact.


\begin{itemize}[leftmargin=1.8em]

\item \textbf{Positive obligation:} the manuscript must perform this role explicitly, not merely avoid a corresponding defect

\item \textbf{Quantitative witness:} the criterion records anchors, counts, coverage booleans, hashes, or repeated-run stability rather than relying on editorial impression alone

\item \textbf{Repair route:} a failure routes to manuscript integration, scientific editorial, proof repair, empirical replay, comparator repair, public metadata repair, or post-release incident handling according to the affected object class

\end{itemize}

\subsection{4. Prior art and scientific context}

A publication-grade manuscript must show that it knows the scientific neighbourhood in which it is speaking. For OC this includes general systems theory, autopoiesis, dynamical systems, hybrid systems, category and type-theoretic formalisms, RAF chemistry, complexity and information measures, identity and persistence theories, systems engineering, reproducible research, artifact evaluation, and scientific publishing norms.


The positive criterion requires overlap-first comparison. The manuscript must state where OC reuses existing ideas, where it reorganizes them, where it adds a residual model-core contribution, and where the available evidence is not strong enough for a priority or unrestricted comparative claim. Literature is not decoration; it is a boundary condition for novelty.


Before writing, the editorial system treats prior art and scientific context as a design requirement rather than as a late-stage defect class. The outline has to allocate a place for the requirement, the corpus register has to identify the supporting sources, and the claim boundary has to name the evidence class that will carry the requirement. This prevents the manuscript from growing by accidental appendix accretion and forces the scientific argument to declare why the material belongs where it appears.


During manuscript integration, prior art and scientific context is checked against three questions. First, does the section teach the reader something that is necessary for judging the released scientific object? Second, does the section cite or explain the artifact layer that supports the statement? Third, does the section give a hostile reviewer a concrete way to challenge the statement? If one of those questions cannot be answered from the text, the section is not ready for release.


At acceptance time, prior art and scientific context must leave quantitative traces. The checker records anchors, section counts, evidence references, reader-question coverage, current-science inclusion, repeated boilerplate ratio, stale-identity hits, and public-surface consistency. The exact numbers are not the scientific claim; they are the guardrail that keeps editorial judgment from becoming a vague impression after the fact.


\begin{itemize}[leftmargin=1.8em]

\item \textbf{Positive obligation:} the manuscript must perform this role explicitly, not merely avoid a corresponding defect

\item \textbf{Quantitative witness:} the criterion records anchors, counts, coverage booleans, hashes, or repeated-run stability rather than relying on editorial impression alone

\item \textbf{Repair route:} a failure routes to manuscript integration, scientific editorial, proof repair, empirical replay, comparator repair, public metadata repair, or post-release incident handling according to the affected object class

\end{itemize}

\subsection{5. Formal proof readability}

Formal sections must not become theorem theatre. A theorem label is useful only when the reader can see the assumptions, definitions, dependency route, proof idea, mechanized subset if available, finite witness if relevant, and counterexample boundary. A proof sheet that only says that a statement follows by definition is not enough for a public theorem claim unless the definition itself carries the whole claimed result and the manuscript says so explicitly.


The positive criterion measures whether theorem names are connected to readable proof surfaces. The reader must be able to move from T133-K0-RES, T133-OMEGA-STATUS, T133-HYBRID, T133-KLEVEL, or T133-MIN to the relevant assumptions, semantic obligations, and falsification routes. A missing route blocks promotion.


Before writing, the editorial system treats formal proof readability as a design requirement rather than as a late-stage defect class. The outline has to allocate a place for the requirement, the corpus register has to identify the supporting sources, and the claim boundary has to name the evidence class that will carry the requirement. This prevents the manuscript from growing by accidental appendix accretion and forces the scientific argument to declare why the material belongs where it appears.


During manuscript integration, formal proof readability is checked against three questions. First, does the section teach the reader something that is necessary for judging the released scientific object? Second, does the section cite or explain the artifact layer that supports the statement? Third, does the section give a hostile reviewer a concrete way to challenge the statement? If one of those questions cannot be answered from the text, the section is not ready for release.


At acceptance time, formal proof readability must leave quantitative traces. The checker records anchors, section counts, evidence references, reader-question coverage, current-science inclusion, repeated boilerplate ratio, stale-identity hits, and public-surface consistency. The exact numbers are not the scientific claim; they are the guardrail that keeps editorial judgment from becoming a vague impression after the fact.


\begin{itemize}[leftmargin=1.8em]

\item \textbf{Positive obligation:} the manuscript must perform this role explicitly, not merely avoid a corresponding defect

\item \textbf{Quantitative witness:} the criterion records anchors, counts, coverage booleans, hashes, or repeated-run stability rather than relying on editorial impression alone

\item \textbf{Repair route:} a failure routes to manuscript integration, scientific editorial, proof repair, empirical replay, comparator repair, public metadata repair, or post-release incident handling according to the affected object class

\end{itemize}

\subsection{6. Empirical and computational evidence}

Empirical language must be stricter than ordinary release language. A promoted empirical row needs a formula or reconstruction rule, an input snapshot, a split or target-blind policy where applicable, a predicted value, an observed value, an uncertainty or residual, a comparator baseline, a negative control, a falsifier, and a replay hash. If the row is only an archive-access test or a known-constant reconstruction, the manuscript must say that and must not upgrade it into full domain evidence-boundary review.


The positive criterion separates numeric replay QA from broad validation. It asks what was predicted or reconstructed, what was observed, how far off it was, what baseline it was compared with, what would count as a negative control, and which result would reopen the claim. Without those fields, empirical promotion is blocked.


Before writing, the editorial system treats empirical and computational evidence as a design requirement rather than as a late-stage defect class. The outline has to allocate a place for the requirement, the corpus register has to identify the supporting sources, and the claim boundary has to name the evidence class that will carry the requirement. This prevents the manuscript from growing by accidental appendix accretion and forces the scientific argument to declare why the material belongs where it appears.


During manuscript integration, empirical and computational evidence is checked against three questions. First, does the section teach the reader something that is necessary for judging the released scientific object? Second, does the section cite or explain the artifact layer that supports the statement? Third, does the section give a hostile reviewer a concrete way to challenge the statement? If one of those questions cannot be answered from the text, the section is not ready for release.


At acceptance time, empirical and computational evidence must leave quantitative traces. The checker records anchors, section counts, evidence references, reader-question coverage, current-science inclusion, repeated boilerplate ratio, stale-identity hits, and public-surface consistency. The exact numbers are not the scientific claim; they are the guardrail that keeps editorial judgment from becoming a vague impression after the fact.


\begin{itemize}[leftmargin=1.8em]

\item \textbf{Positive obligation:} the manuscript must perform this role explicitly, not merely avoid a corresponding defect

\item \textbf{Quantitative witness:} the criterion records anchors, counts, coverage booleans, hashes, or repeated-run stability rather than relying on editorial impression alone

\item \textbf{Repair route:} a failure routes to manuscript integration, scientific editorial, proof repair, empirical replay, comparator repair, public metadata repair, or post-release incident handling according to the affected object class

\end{itemize}

\subsection{7. Phenomenon coverage}

A model may be ambitious only if each phenomenon route is honest about what it covers. A phenomenon card must specify the OC instance, the observable, the explanation route, the comparator, the negative control, the falsifier, and the current status. The manuscript must distinguish explained, protocol-ready, partially covered, blocked, and future-research phenomena.


The positive criterion rejects coverage inflation. Naming a phenomenon is not evidence that the phenomenon has been explained. The text must show how the model reaches it, what remains outside the release, and what a hostile reviewer could test.


Before writing, the editorial system treats phenomenon coverage as a design requirement rather than as a late-stage defect class. The outline has to allocate a place for the requirement, the corpus register has to identify the supporting sources, and the claim boundary has to name the evidence class that will carry the requirement. This prevents the manuscript from growing by accidental appendix accretion and forces the scientific argument to declare why the material belongs where it appears.


During manuscript integration, phenomenon coverage is checked against three questions. First, does the section teach the reader something that is necessary for judging the released scientific object? Second, does the section cite or explain the artifact layer that supports the statement? Third, does the section give a hostile reviewer a concrete way to challenge the statement? If one of those questions cannot be answered from the text, the section is not ready for release.


At acceptance time, phenomenon coverage must leave quantitative traces. The checker records anchors, section counts, evidence references, reader-question coverage, current-science inclusion, repeated boilerplate ratio, stale-identity hits, and public-surface consistency. The exact numbers are not the scientific claim; they are the guardrail that keeps editorial judgment from becoming a vague impression after the fact.


\begin{itemize}[leftmargin=1.8em]

\item \textbf{Positive obligation:} the manuscript must perform this role explicitly, not merely avoid a corresponding defect

\item \textbf{Quantitative witness:} the criterion records anchors, counts, coverage booleans, hashes, or repeated-run stability rather than relying on editorial impression alone

\item \textbf{Repair route:} a failure routes to manuscript integration, scientific editorial, proof repair, empirical replay, comparator repair, public metadata repair, or post-release incident handling according to the affected object class

\end{itemize}

\subsection{8. Novelty and equivalence pressure}

The strongest novelty attack says that OC is another theory described from a different angle. A publication-grade release must answer that attack without slogans. The answer has to compare specific claims against specific traditions, state overlap, identify residual delta, and explain why the residual delta matters for the model core. Where the residual delta is not proved, the novelty claim must be bounded.


The positive criterion checks that novelty rows have named comparators, overlap fields, residual-delta fields, and uniqueness status. It also checks that the public manuscript does not convert bounded residual delta into scoped priority or total scientific unrestricted comparative claim.


Before writing, the editorial system treats novelty and equivalence pressure as a design requirement rather than as a late-stage defect class. The outline has to allocate a place for the requirement, the corpus register has to identify the supporting sources, and the claim boundary has to name the evidence class that will carry the requirement. This prevents the manuscript from growing by accidental appendix accretion and forces the scientific argument to declare why the material belongs where it appears.


During manuscript integration, novelty and equivalence pressure is checked against three questions. First, does the section teach the reader something that is necessary for judging the released scientific object? Second, does the section cite or explain the artifact layer that supports the statement? Third, does the section give a hostile reviewer a concrete way to challenge the statement? If one of those questions cannot be answered from the text, the section is not ready for release.


At acceptance time, novelty and equivalence pressure must leave quantitative traces. The checker records anchors, section counts, evidence references, reader-question coverage, current-science inclusion, repeated boilerplate ratio, stale-identity hits, and public-surface consistency. The exact numbers are not the scientific claim; they are the guardrail that keeps editorial judgment from becoming a vague impression after the fact.


\begin{itemize}[leftmargin=1.8em]

\item \textbf{Positive obligation:} the manuscript must perform this role explicitly, not merely avoid a corresponding defect

\item \textbf{Quantitative witness:} the criterion records anchors, counts, coverage booleans, hashes, or repeated-run stability rather than relying on editorial impression alone

\item \textbf{Repair route:} a failure routes to manuscript integration, scientific editorial, proof repair, empirical replay, comparator repair, public metadata repair, or post-release incident handling according to the affected object class

\end{itemize}

\subsection{9. Didactic progression and figures}

A public scientific text should teach before it exhausts. The reader should encounter a route from tuple to liveness, from liveness to boundary, from boundary to operator, from operator to K-level witness, from K-level witness to proof or replay, and from proof or replay to falsifier. Figures, route tables, and worked examples are therefore part of the scientific argument, not cosmetic features.


The positive criterion asks whether the text has visual pedagogy anchors, figure routes, worked examples, and transitions. A document with many registers but no teachable path fails because it forces the reader to reconstruct the theory from storage layout rather than from scientific exposition.


Before writing, the editorial system treats didactic progression and figures as a design requirement rather than as a late-stage defect class. The outline has to allocate a place for the requirement, the corpus register has to identify the supporting sources, and the claim boundary has to name the evidence class that will carry the requirement. This prevents the manuscript from growing by accidental appendix accretion and forces the scientific argument to declare why the material belongs where it appears.


During manuscript integration, didactic progression and figures is checked against three questions. First, does the section teach the reader something that is necessary for judging the released scientific object? Second, does the section cite or explain the artifact layer that supports the statement? Third, does the section give a hostile reviewer a concrete way to challenge the statement? If one of those questions cannot be answered from the text, the section is not ready for release.


At acceptance time, didactic progression and figures must leave quantitative traces. The checker records anchors, section counts, evidence references, reader-question coverage, current-science inclusion, repeated boilerplate ratio, stale-identity hits, and public-surface consistency. The exact numbers are not the scientific claim; they are the guardrail that keeps editorial judgment from becoming a vague impression after the fact.


\begin{itemize}[leftmargin=1.8em]

\item \textbf{Positive obligation:} the manuscript must perform this role explicitly, not merely avoid a corresponding defect

\item \textbf{Quantitative witness:} the criterion records anchors, counts, coverage booleans, hashes, or repeated-run stability rather than relying on editorial impression alone

\item \textbf{Repair route:} a failure routes to manuscript integration, scientific editorial, proof repair, empirical replay, comparator repair, public metadata repair, or post-release incident handling according to the affected object class

\end{itemize}

\subsection{10. Editorial unity and style}

A monograph is one work, not a stack of unrelated artifacts. Page numbering, section order, terminology, title pages, abstracts, dedication, citations, and conclusion must behave as a single manuscript. Support documents must share the same object model and claim boundary. The result should read as an edited scientific work, not as a repository export.


The positive criterion checks for coherent front matter, dedication where appropriate, table of contents, section transitions, repeated-boilerplate ratio, raw TeX leakage, raw JSON leakage, checksum-only dumps, status dumps, stale versions, and absent conclusion. The editorial target is not prettiness alone; style is how the scientific object remains intelligible.


Before writing, the editorial system treats editorial unity and style as a design requirement rather than as a late-stage defect class. The outline has to allocate a place for the requirement, the corpus register has to identify the supporting sources, and the claim boundary has to name the evidence class that will carry the requirement. This prevents the manuscript from growing by accidental appendix accretion and forces the scientific argument to declare why the material belongs where it appears.


During manuscript integration, editorial unity and style is checked against three questions. First, does the section teach the reader something that is necessary for judging the released scientific object? Second, does the section cite or explain the artifact layer that supports the statement? Third, does the section give a hostile reviewer a concrete way to challenge the statement? If one of those questions cannot be answered from the text, the section is not ready for release.


At acceptance time, editorial unity and style must leave quantitative traces. The checker records anchors, section counts, evidence references, reader-question coverage, current-science inclusion, repeated boilerplate ratio, stale-identity hits, and public-surface consistency. The exact numbers are not the scientific claim; they are the guardrail that keeps editorial judgment from becoming a vague impression after the fact.


\begin{itemize}[leftmargin=1.8em]

\item \textbf{Positive obligation:} the manuscript must perform this role explicitly, not merely avoid a corresponding defect

\item \textbf{Quantitative witness:} the criterion records anchors, counts, coverage booleans, hashes, or repeated-run stability rather than relying on editorial impression alone

\item \textbf{Repair route:} a failure routes to manuscript integration, scientific editorial, proof repair, empirical replay, comparator repair, public metadata repair, or post-release incident handling according to the affected object class

\end{itemize}

\subsection{11. Internal ontology and external positioning}

The release must distinguish persons, instruments, methods, corpora, projects, workstreams, release records, journal packages, evidence artifacts, and governance states. Alexander Yashin is the author; the release engineering pipeline is the research-instrument and institute-automation system; the methodological staging framework is the methodological framework; OC Core is the scientific model release. Internal object names do not automatically become external author, affiliation, or claim labels.


The positive criterion checks public metadata and document language against the entity ontology. Misclassifying an instrument as an author, a method as an organization, an internal status as a public claim, or a journal owner-review packet as a submitted article is a publication failure.


Before writing, the editorial system treats internal ontology and external positioning as a design requirement rather than as a late-stage defect class. The outline has to allocate a place for the requirement, the corpus register has to identify the supporting sources, and the claim boundary has to name the evidence class that will carry the requirement. This prevents the manuscript from growing by accidental appendix accretion and forces the scientific argument to declare why the material belongs where it appears.


During manuscript integration, internal ontology and external positioning is checked against three questions. First, does the section teach the reader something that is necessary for judging the released scientific object? Second, does the section cite or explain the artifact layer that supports the statement? Third, does the section give a hostile reviewer a concrete way to challenge the statement? If one of those questions cannot be answered from the text, the section is not ready for release.


At acceptance time, internal ontology and external positioning must leave quantitative traces. The checker records anchors, section counts, evidence references, reader-question coverage, current-science inclusion, repeated boilerplate ratio, stale-identity hits, and public-surface consistency. The exact numbers are not the scientific claim; they are the guardrail that keeps editorial judgment from becoming a vague impression after the fact.


\begin{itemize}[leftmargin=1.8em]

\item \textbf{Positive obligation:} the manuscript must perform this role explicitly, not merely avoid a corresponding defect

\item \textbf{Quantitative witness:} the criterion records anchors, counts, coverage booleans, hashes, or repeated-run stability rather than relying on editorial impression alone

\item \textbf{Repair route:} a failure routes to manuscript integration, scientific editorial, proof repair, empirical replay, comparator repair, public metadata repair, or post-release incident handling according to the affected object class

\end{itemize}

\subsection{12. Release packaging and archive presentation}

A scientific archive record is part of the publication surface. Its title, abstract, description, files, creators, ORCID, license, keywords, related identifiers, DOI, concept DOI, reading order, and file preview must describe the same scientific object as the PDFs. Archive metadata must not be a checksum-only dump, unrendered platform markup, or process-governance dump.


The positive criterion therefore checks public metadata as text. It verifies that GitHub and Zenodo have distinct renderers, matching DOI and asset sets, professional descriptions, and curated public files. Journal packages may be included as owner-review material, but they must not imply that submission has occurred.


Before writing, the editorial system treats release packaging and archive presentation as a design requirement rather than as a late-stage defect class. The outline has to allocate a place for the requirement, the corpus register has to identify the supporting sources, and the claim boundary has to name the evidence class that will carry the requirement. This prevents the manuscript from growing by accidental appendix accretion and forces the scientific argument to declare why the material belongs where it appears.


During manuscript integration, release packaging and archive presentation is checked against three questions. First, does the section teach the reader something that is necessary for judging the released scientific object? Second, does the section cite or explain the artifact layer that supports the statement? Third, does the section give a hostile reviewer a concrete way to challenge the statement? If one of those questions cannot be answered from the text, the section is not ready for release.


At acceptance time, release packaging and archive presentation must leave quantitative traces. The checker records anchors, section counts, evidence references, reader-question coverage, current-science inclusion, repeated boilerplate ratio, stale-identity hits, and public-surface consistency. The exact numbers are not the scientific claim; they are the guardrail that keeps editorial judgment from becoming a vague impression after the fact.


\begin{itemize}[leftmargin=1.8em]

\item \textbf{Positive obligation:} the manuscript must perform this role explicitly, not merely avoid a corresponding defect

\item \textbf{Quantitative witness:} the criterion records anchors, counts, coverage booleans, hashes, or repeated-run stability rather than relying on editorial impression alone

\item \textbf{Repair route:} a failure routes to manuscript integration, scientific editorial, proof repair, empirical replay, comparator repair, public metadata repair, or post-release incident handling according to the affected object class

\end{itemize}

\subsection{13. Reproducibility and delta discipline}

Verification should not damage the release it verifies. A read-only check that rewrites artifacts without semantic change is a process defect. A meaningful semantic delta should trigger the minimum necessary regeneration and review route; absence of delta should stop the downstream chain. This prevents the machine from fighting itself and lets editorial attention go to real changes.


The positive criterion records repeat stability, checksum consistency, dirty-tree governance, and current-science-state summary state. It is not enough to pass once after regenerating everything; the pipeline must show that repeated checks leave the artifact stable when inputs are unchanged.


Before writing, the editorial system treats reproducibility and delta discipline as a design requirement rather than as a late-stage defect class. The outline has to allocate a place for the requirement, the corpus register has to identify the supporting sources, and the claim boundary has to name the evidence class that will carry the requirement. This prevents the manuscript from growing by accidental appendix accretion and forces the scientific argument to declare why the material belongs where it appears.


During manuscript integration, reproducibility and delta discipline is checked against three questions. First, does the section teach the reader something that is necessary for judging the released scientific object? Second, does the section cite or explain the artifact layer that supports the statement? Third, does the section give a hostile reviewer a concrete way to challenge the statement? If one of those questions cannot be answered from the text, the section is not ready for release.


At acceptance time, reproducibility and delta discipline must leave quantitative traces. The checker records anchors, section counts, evidence references, reader-question coverage, current-science inclusion, repeated boilerplate ratio, stale-identity hits, and public-surface consistency. The exact numbers are not the scientific claim; they are the guardrail that keeps editorial judgment from becoming a vague impression after the fact.


\begin{itemize}[leftmargin=1.8em]

\item \textbf{Positive obligation:} the manuscript must perform this role explicitly, not merely avoid a corresponding defect

\item \textbf{Quantitative witness:} the criterion records anchors, counts, coverage booleans, hashes, or repeated-run stability rather than relying on editorial impression alone

\item \textbf{Repair route:} a failure routes to manuscript integration, scientific editorial, proof repair, empirical replay, comparator repair, public metadata repair, or post-release incident handling according to the affected object class

\end{itemize}

\subsection{14. Known-error learning}

A failed publication must become a permanent test. The 1.3.3 incident produced known-error classes: process memos in scientific roles, unrendered platform markup in archive descriptions, tiny surrogate PDFs, absent title pages, absent dedication, stale version identity, archive records with metadata-first file preview, appended deltas, page-number discontinuity, unsupported public claims, and generated status rows masquerading as prose.


The positive criterion is paired with a known-error gate. Negative patterns block recurrence, while positive obligations demand replacement with a real scientific artifact. Learning from the incident means future releases fail before publication when the same behaviour appears, even if the file names look correct.


Before writing, the editorial system treats known-error learning as a design requirement rather than as a late-stage defect class. The outline has to allocate a place for the requirement, the corpus register has to identify the supporting sources, and the claim boundary has to name the evidence class that will carry the requirement. This prevents the manuscript from growing by accidental appendix accretion and forces the scientific argument to declare why the material belongs where it appears.


During manuscript integration, known-error learning is checked against three questions. First, does the section teach the reader something that is necessary for judging the released scientific object? Second, does the section cite or explain the artifact layer that supports the statement? Third, does the section give a hostile reviewer a concrete way to challenge the statement? If one of those questions cannot be answered from the text, the section is not ready for release.


At acceptance time, known-error learning must leave quantitative traces. The checker records anchors, section counts, evidence references, reader-question coverage, current-science inclusion, repeated boilerplate ratio, stale-identity hits, and public-surface consistency. The exact numbers are not the scientific claim; they are the guardrail that keeps editorial judgment from becoming a vague impression after the fact.


\begin{itemize}[leftmargin=1.8em]

\item \textbf{Positive obligation:} the manuscript must perform this role explicitly, not merely avoid a corresponding defect

\item \textbf{Quantitative witness:} the criterion records anchors, counts, coverage booleans, hashes, or repeated-run stability rather than relying on editorial impression alone

\item \textbf{Repair route:} a failure routes to manuscript integration, scientific editorial, proof repair, empirical replay, comparator repair, public metadata repair, or post-release incident handling according to the affected object class

\end{itemize}

\subsection{15. Post-release accountability}

Publication does not end the scientific obligation. A post-release check must fetch public records, verify titles, descriptions, files, checksums, DOI, tag target, downloadable PDFs, archive rendering, citation metadata, and claim boundaries. If the public surface differs from the local verified surface, the release is reopened as a publication incident.


The positive criterion treats public reality as the final artifact. Local success is necessary but not sufficient. The public record must be readable, citable, coherent, and scientifically bounded after the platform has rendered it.


Before writing, the editorial system treats post-release accountability as a design requirement rather than as a late-stage defect class. The outline has to allocate a place for the requirement, the corpus register has to identify the supporting sources, and the claim boundary has to name the evidence class that will carry the requirement. This prevents the manuscript from growing by accidental appendix accretion and forces the scientific argument to declare why the material belongs where it appears.


During manuscript integration, post-release accountability is checked against three questions. First, does the section teach the reader something that is necessary for judging the released scientific object? Second, does the section cite or explain the artifact layer that supports the statement? Third, does the section give a hostile reviewer a concrete way to challenge the statement? If one of those questions cannot be answered from the text, the section is not ready for release.


At acceptance time, post-release accountability must leave quantitative traces. The checker records anchors, section counts, evidence references, reader-question coverage, current-science inclusion, repeated boilerplate ratio, stale-identity hits, and public-surface consistency. The exact numbers are not the scientific claim; they are the guardrail that keeps editorial judgment from becoming a vague impression after the fact.


\begin{itemize}[leftmargin=1.8em]

\item \textbf{Positive obligation:} the manuscript must perform this role explicitly, not merely avoid a corresponding defect

\item \textbf{Quantitative witness:} the criterion records anchors, counts, coverage booleans, hashes, or repeated-run stability rather than relying on editorial impression alone

\item \textbf{Repair route:} a failure routes to manuscript integration, scientific editorial, proof repair, empirical replay, comparator repair, public metadata repair, or post-release incident handling according to the affected object class

\end{itemize}

\section{Scientific Editorial Board Protocol}

\label{sec:oc133-scientific-editorial-board-protocol}

A release reaches publication quality only after the manuscript has been read through multiple editorial roles. The roles below are not honorary labels. Each role has a concrete defect class, a positive must-have class, and a routing consequence. A role may approve only the object class it is competent to judge; no role may convert a missing proof, missing data row, weak public description, or broken layout into acceptance by stylistic approval alone.


The board protocol is part of the publication standard because complex scientific artifacts fail at interfaces. A theorem can be formally careful but unreadable; an archive can contain correct files but present them badly; an empirical table can be reproducible but over-advertised; a manuscript can have good sections but no narrative route. The protocol makes those interfaces reviewable before publication.


\subsection{Editorial role 1: Scientific acquisitions editor}

The scientific acquisitions editor asks whether the manuscript has a real scientific reason to exist, whether the contribution is framed at the right scale, and whether the release object is the correct channel for the claim surface. This role reads the artifact as a specialist with authority over a particular criticism route rather than as a general-purpose approval stamp.


Acceptance by this role requires a clear thesis, a bounded contribution statement, a reader contract, a target audience, and a statement of what would make the release scientifically premature. If the role cannot produce a clear pass/fail basis, the manuscript returns to the relevant service: research, proof repair, empirical replay, comparator work, manuscript integration, metadata repair, or incident management.


The role also leaves a reusable trace. The trace states the inspected artifact, the requirement, the evidence used for acceptance, the defect that would reopen the finding, and the regression test that prevents the same class of error from recurring in a later release.


\begin{itemize}[leftmargin=1.8em]

\item \textbf{Primary object:} public manuscript, support PDF, evidence package, metadata record, or post-release public surface

\item \textbf{Positive acceptance:} the object fulfills its scientific and editorial purpose in full, not merely avoids an explicit defect

\item \textbf{Reopen condition:} a missing route, weak explanation, unsupported claim, broken evidence binding, public rendering mismatch, or known-error recurrence

\end{itemize}

\subsection{Editorial role 2: Formal proof editor}

The formal proof editor checks theorem labels, assumptions, definitions, dependency routes, proof sheets, mechanized subset declarations, finite witness boundaries, and counterexample conditions. This role reads the artifact as a specialist with authority over a particular criticism route rather than as a general-purpose approval stamp.


Acceptance by this role requires each promoted theorem-like statement to have a visible proof route or to be demoted before public release. If the role cannot produce a clear pass/fail basis, the manuscript returns to the relevant service: research, proof repair, empirical replay, comparator work, manuscript integration, metadata repair, or incident management.


The role also leaves a reusable trace. The trace states the inspected artifact, the requirement, the evidence used for acceptance, the defect that would reopen the finding, and the regression test that prevents the same class of error from recurring in a later release.


\begin{itemize}[leftmargin=1.8em]

\item \textbf{Primary object:} public manuscript, support PDF, evidence package, metadata record, or post-release public surface

\item \textbf{Positive acceptance:} the object fulfills its scientific and editorial purpose in full, not merely avoids an explicit defect

\item \textbf{Reopen condition:} a missing route, weak explanation, unsupported claim, broken evidence binding, public rendering mismatch, or known-error recurrence

\end{itemize}

\subsection{Editorial role 3: Empirical statistics editor}

The empirical statistics editor checks formulas, data snapshots, target-blind or held-out policy, predicted and observed values, residuals, uncertainties, comparators, negative controls, falsifiers, and replay hashes. This role reads the artifact as a specialist with authority over a particular criticism route rather than as a general-purpose approval stamp.


Acceptance by this role requires empirical promotion to be impossible without numeric or executable evidence matching the public wording. If the role cannot produce a clear pass/fail basis, the manuscript returns to the relevant service: research, proof repair, empirical replay, comparator work, manuscript integration, metadata repair, or incident management.


The role also leaves a reusable trace. The trace states the inspected artifact, the requirement, the evidence used for acceptance, the defect that would reopen the finding, and the regression test that prevents the same class of error from recurring in a later release.


\begin{itemize}[leftmargin=1.8em]

\item \textbf{Primary object:} public manuscript, support PDF, evidence package, metadata record, or post-release public surface

\item \textbf{Positive acceptance:} the object fulfills its scientific and editorial purpose in full, not merely avoids an explicit defect

\item \textbf{Reopen condition:} a missing route, weak explanation, unsupported claim, broken evidence binding, public rendering mismatch, or known-error recurrence

\end{itemize}

\subsection{Editorial role 4: Prior-art and novelty editor}

The prior-art and novelty editor checks whether the manuscript fairly represents adjacent traditions, whether overlap is stated before residual delta, and whether priority language is supported by a source-backed comparator row. This role reads the artifact as a specialist with authority over a particular criticism route rather than as a general-purpose approval stamp.


Acceptance by this role requires unsupported novelty, equivalence, unrestricted comparative claim, and total-synthesis language to be removed or moved to outside-promoted-surface status. If the role cannot produce a clear pass/fail basis, the manuscript returns to the relevant service: research, proof repair, empirical replay, comparator work, manuscript integration, metadata repair, or incident management.


The role also leaves a reusable trace. The trace states the inspected artifact, the requirement, the evidence used for acceptance, the defect that would reopen the finding, and the regression test that prevents the same class of error from recurring in a later release.


\begin{itemize}[leftmargin=1.8em]

\item \textbf{Primary object:} public manuscript, support PDF, evidence package, metadata record, or post-release public surface

\item \textbf{Positive acceptance:} the object fulfills its scientific and editorial purpose in full, not merely avoids an explicit defect

\item \textbf{Reopen condition:} a missing route, weak explanation, unsupported claim, broken evidence binding, public rendering mismatch, or known-error recurrence

\end{itemize}

\subsection{Editorial role 5: Phenomenon coverage editor}

The phenomenon coverage editor checks whether every named phenomenon has an OC instance, observable, explanation route, comparator, negative control, falsifier, and honest status. This role reads the artifact as a specialist with authority over a particular criticism route rather than as a general-purpose approval stamp.


Acceptance by this role requires phenomenon lists to distinguish explained, protocol-ready, partial, blocked, and future-research cases. If the role cannot produce a clear pass/fail basis, the manuscript returns to the relevant service: research, proof repair, empirical replay, comparator work, manuscript integration, metadata repair, or incident management.


The role also leaves a reusable trace. The trace states the inspected artifact, the requirement, the evidence used for acceptance, the defect that would reopen the finding, and the regression test that prevents the same class of error from recurring in a later release.


\begin{itemize}[leftmargin=1.8em]

\item \textbf{Primary object:} public manuscript, support PDF, evidence package, metadata record, or post-release public surface

\item \textbf{Positive acceptance:} the object fulfills its scientific and editorial purpose in full, not merely avoids an explicit defect

\item \textbf{Reopen condition:} a missing route, weak explanation, unsupported claim, broken evidence binding, public rendering mismatch, or known-error recurrence

\end{itemize}

\subsection{Editorial role 6: Didactic editor}

The didactic editor checks whether a prepared but hostile reader can learn the tuple, liveness/death/residue distinction, K-level route, boundary classifier, operator semantics, evidence route, and claim limits in order. This role reads the artifact as a specialist with authority over a particular criticism route rather than as a general-purpose approval stamp.


Acceptance by this role requires transitions, worked examples, visual anchors, and explanations before large registers or appendices can carry interpretive weight. If the role cannot produce a clear pass/fail basis, the manuscript returns to the relevant service: research, proof repair, empirical replay, comparator work, manuscript integration, metadata repair, or incident management.


The role also leaves a reusable trace. The trace states the inspected artifact, the requirement, the evidence used for acceptance, the defect that would reopen the finding, and the regression test that prevents the same class of error from recurring in a later release.


\begin{itemize}[leftmargin=1.8em]

\item \textbf{Primary object:} public manuscript, support PDF, evidence package, metadata record, or post-release public surface

\item \textbf{Positive acceptance:} the object fulfills its scientific and editorial purpose in full, not merely avoids an explicit defect

\item \textbf{Reopen condition:} a missing route, weak explanation, unsupported claim, broken evidence binding, public rendering mismatch, or known-error recurrence

\end{itemize}

\subsection{Editorial role 7: Copyeditor}

The copyeditor checks sentence shape, terminology consistency, repeated boilerplate, undefined abbreviations, punctuation, stale references, and awkward machine-generated phrasing. This role reads the artifact as a specialist with authority over a particular criticism route rather than as a general-purpose approval stamp.


Acceptance by this role requires the manuscript to read as a unified scientific work rather than as stitched automation output. If the role cannot produce a clear pass/fail basis, the manuscript returns to the relevant service: research, proof repair, empirical replay, comparator work, manuscript integration, metadata repair, or incident management.


The role also leaves a reusable trace. The trace states the inspected artifact, the requirement, the evidence used for acceptance, the defect that would reopen the finding, and the regression test that prevents the same class of error from recurring in a later release.


\begin{itemize}[leftmargin=1.8em]

\item \textbf{Primary object:} public manuscript, support PDF, evidence package, metadata record, or post-release public surface

\item \textbf{Positive acceptance:} the object fulfills its scientific and editorial purpose in full, not merely avoids an explicit defect

\item \textbf{Reopen condition:} a missing route, weak explanation, unsupported claim, broken evidence binding, public rendering mismatch, or known-error recurrence

\end{itemize}

\subsection{Editorial role 8: Layout and production editor}

The layout and production editor checks title pages, dedication, abstract, table of contents, page numbering, figure and table flow, PDF metadata, readable tables, line breaks, and archive-preview order. This role reads the artifact as a specialist with authority over a particular criticism route rather than as a general-purpose approval stamp.


Acceptance by this role requires every public PDF to be self-identifying, visually coherent, and usable as a standalone scientific artifact. If the role cannot produce a clear pass/fail basis, the manuscript returns to the relevant service: research, proof repair, empirical replay, comparator work, manuscript integration, metadata repair, or incident management.


The role also leaves a reusable trace. The trace states the inspected artifact, the requirement, the evidence used for acceptance, the defect that would reopen the finding, and the regression test that prevents the same class of error from recurring in a later release.


\begin{itemize}[leftmargin=1.8em]

\item \textbf{Primary object:} public manuscript, support PDF, evidence package, metadata record, or post-release public surface

\item \textbf{Positive acceptance:} the object fulfills its scientific and editorial purpose in full, not merely avoids an explicit defect

\item \textbf{Reopen condition:} a missing route, weak explanation, unsupported claim, broken evidence binding, public rendering mismatch, or known-error recurrence

\end{itemize}

\subsection{Editorial role 9: Bibliography and source editor}

The bibliography and source editor checks whether the reference list represents the scientific neighbourhood of the work and whether cited claims can be located in the cited traditions or in the release evidence package. This role reads the artifact as a specialist with authority over a particular criticism route rather than as a general-purpose approval stamp.


Acceptance by this role requires the bibliography to support context, novelty boundaries, and reproducibility rather than serving as a decorative list. If the role cannot produce a clear pass/fail basis, the manuscript returns to the relevant service: research, proof repair, empirical replay, comparator work, manuscript integration, metadata repair, or incident management.


The role also leaves a reusable trace. The trace states the inspected artifact, the requirement, the evidence used for acceptance, the defect that would reopen the finding, and the regression test that prevents the same class of error from recurring in a later release.


\begin{itemize}[leftmargin=1.8em]

\item \textbf{Primary object:} public manuscript, support PDF, evidence package, metadata record, or post-release public surface

\item \textbf{Positive acceptance:} the object fulfills its scientific and editorial purpose in full, not merely avoids an explicit defect

\item \textbf{Reopen condition:} a missing route, weak explanation, unsupported claim, broken evidence binding, public rendering mismatch, or known-error recurrence

\end{itemize}

\subsection{Editorial role 10: Public metadata editor}

The public metadata editor checks archive title, abstract, creators, ORCID, license, keywords, DOI, concept DOI, related identifiers, release body, file set, checksums, and rendered description. This role reads the artifact as a specialist with authority over a particular criticism route rather than as a general-purpose approval stamp.


Acceptance by this role requires public platforms to display a professional scientific landing surface instead of a raw control dump. If the role cannot produce a clear pass/fail basis, the manuscript returns to the relevant service: research, proof repair, empirical replay, comparator work, manuscript integration, metadata repair, or incident management.


The role also leaves a reusable trace. The trace states the inspected artifact, the requirement, the evidence used for acceptance, the defect that would reopen the finding, and the regression test that prevents the same class of error from recurring in a later release.


\begin{itemize}[leftmargin=1.8em]

\item \textbf{Primary object:} public manuscript, support PDF, evidence package, metadata record, or post-release public surface

\item \textbf{Positive acceptance:} the object fulfills its scientific and editorial purpose in full, not merely avoids an explicit defect

\item \textbf{Reopen condition:} a missing route, weak explanation, unsupported claim, broken evidence binding, public rendering mismatch, or known-error recurrence

\end{itemize}

\subsection{Editorial role 11: Incident and known-error editor}

The incident and known-error editor checks whether every previous criticism route has become a test, whether the repair addresses the behaviour class, and whether post-release verification can detect platform-side drift. This role reads the artifact as a specialist with authority over a particular criticism route rather than as a general-purpose approval stamp.


Acceptance by this role requires the machine to learn from publication defects before a corrected release is allowed to replace the public record. If the role cannot produce a clear pass/fail basis, the manuscript returns to the relevant service: research, proof repair, empirical replay, comparator work, manuscript integration, metadata repair, or incident management.


The role also leaves a reusable trace. The trace states the inspected artifact, the requirement, the evidence used for acceptance, the defect that would reopen the finding, and the regression test that prevents the same class of error from recurring in a later release.


\begin{itemize}[leftmargin=1.8em]

\item \textbf{Primary object:} public manuscript, support PDF, evidence package, metadata record, or post-release public surface

\item \textbf{Positive acceptance:} the object fulfills its scientific and editorial purpose in full, not merely avoids an explicit defect

\item \textbf{Reopen condition:} a missing route, weak explanation, unsupported claim, broken evidence binding, public rendering mismatch, or known-error recurrence

\end{itemize}

\subsection{Editorial role 12: Owner-review editor}

The owner-review editor checks the final package from the standpoint of the human owner: whether the work is scientifically defensible, visually acceptable, citable, properly bounded, and ready for public conversation. This role reads the artifact as a specialist with authority over a particular criticism route rather than as a general-purpose approval stamp.


Acceptance by this role requires unresolved editorial, scientific, metadata, or archive-surface concerns to block publication even when machine checks are otherwise green. If the role cannot produce a clear pass/fail basis, the manuscript returns to the relevant service: research, proof repair, empirical replay, comparator work, manuscript integration, metadata repair, or incident management.


The role also leaves a reusable trace. The trace states the inspected artifact, the requirement, the evidence used for acceptance, the defect that would reopen the finding, and the regression test that prevents the same class of error from recurring in a later release.


\begin{itemize}[leftmargin=1.8em]

\item \textbf{Primary object:} public manuscript, support PDF, evidence package, metadata record, or post-release public surface

\item \textbf{Positive acceptance:} the object fulfills its scientific and editorial purpose in full, not merely avoids an explicit defect

\item \textbf{Reopen condition:} a missing route, weak explanation, unsupported claim, broken evidence binding, public rendering mismatch, or known-error recurrence

\end{itemize}

\subsection{How the Standard Applies to OC Core 1.3.3}

For OC Core 1.3.3 the standard requires the public monograph to carry the model core, proof route, finite and Lean evidence, bounded replay rows, prior-art map, phenomenon coverage, reviewer pressure, reproducibility route, journal owner-review map, and public archive boundary in one coherent manuscript. The supporting PDFs then become entry points into the same scientific object rather than independent process packets.


The standard also states what the release must not do. It must not claim full full-scope completion, a complete final unification claim, or a global comparison victory over contemporary science unless the artifact layer literally supports such claims. It may preserve those ambitions as non-promoted research obligations, but the public release surface must remain evidence-bound.


This appendix is part of the release because the method of publication is itself now an object under review. A reader should be able to judge not only whether the model is interesting, but whether the release machine has learned how to turn current science into readable, bounded, auditable public science.


\subsection{Minimum Release-Button Interpretation}

The practical consequence of this standard is that the release button is not a formatting action. It is a scientific commitment that the current object has been written, checked, reviewed, archived, and bounded as one coherent publication. The button may be pressed only after the text, evidence, metadata, and public presentation describe the same object. A mismatch between those surfaces is not a small production issue; it is a scientific communication failure because it changes what the outside reader is asked to believe.


For OC Core 1.3.3 the release-button interpretation is deliberately bounded. The release may say that the model core has been presented in a publication-grade form with typed foundations, proof routes, selected mechanized evidence, finite semantic witnesses, bounded replay rows, prior-art positioning, phenomenon routes, and adversarial review. It may not say that every possible domain projection has been completed or that every contemporary scientific model has been defeated by a scoped numerical predictor.


The distinction is not a retreat from ambition. It is the condition that lets ambition remain scientifically usable. A large theory becomes reviewable only when the reader can separate the supported present surface from the research program that follows. The supported surface must be strong, explicit, and citable. The research program must be visible, falsifiable, and routed. Confusing the two would make the manuscript louder but weaker.


The positive criterion therefore treats completeness as a routed relation. Completeness does not mean that the text contains every file in the repository. It means that every promoted public claim has the required reader explanation and evidence path, that every current science surface needed for the claim has been represented, and that every excluded or background obligation is named honestly. This is the standard that prevents another release from looking complete in storage while failing as a scientific work.


The standard is intentionally reusable beyond OC. Any the release engineering pipeline-produced scientific publication should be able to inherit the same object discipline: entity ontology first, research state second, manuscript integration third, editorial review fourth, release packaging fifth, public presentation sixth, and post-release verification seventh. The product changes; the publication method remains a governed service.


This ordering matters because each stage protects a different part of the truth condition. Entity ontology protects who and what the objects are. Research state protects what the work currently knows. Manuscript integration protects how knowledge becomes teachable prose. Editorial review protects the reader from unclear, inflated, or malformed communication. Release packaging protects the archived object. Public presentation protects the external first encounter. Post-release verification protects against platform rendering and file drift.


The same ordering also prevents local fixes from becoming future regressions. If an archive description is ugly, the repair is not just to edit the description; the repair is to separate Markdown and HTML renderers, test both, and verify the public page. If a monograph is an appended delta, the repair is not just to merge PDFs differently; the repair is to require source-to-manuscript integration and coverage proof. If a claim is overpromoted, the repair is not just to delete a sentence; the repair is to route claim strength through the evidence class and comparator boundary.


The release-button state is therefore a compositional state. It is obtained only when document quality, scientific support, editorial acceptance, archive metadata, and public surface verification agree. This is why a single green machine criterion is not sufficient after an incident. The corrected machine needs multiple independent positive criteria that converge on the same object from different disciplines.


The expected external result is modest in wording and strong in construction. A scientist opening the release should see a real monograph, a real article-style entry point, a real methods companion, a real reviewer map, a real evidence package, professional metadata, and a coherent relationship among them. The reader should not have to know the internal pipeline to understand what the release claims or how to audit it.


That is the publication-grade target used here. It is stricter than file completeness, stricter than metadata validity, stricter than absence of known bad phrases, and stricter than local reproducibility. It requires the scientific object to be understandable, defensible, inspectable, bounded, and professionally presented as one artifact family.
